\section{Conclusion}
\label{sec:conclusion}

We presented \textbf{RoboDojo}, a unified sim-and-real benchmark for evaluating generalist robot manipulation policies. RoboDojo includes 42 simulation tasks and 18 real-world tasks, covering diverse, long-horizon, precision-demanding, memory-dependent, and open-ended manipulation scenarios. The simulation benchmark supports efficient large-scale evaluation through heterogeneous parallel execution, enabling rapid feedback and fine-grained capability diagnosis. The real-world benchmark, built on RoboDojo-RealEval, provides standardized hardware configurations, scene reset procedures, evaluation protocols, and remote cloud-based access for reproducible physical assessment. Together with \textbf{XPolicyLab}, a unified standard and infrastructure for policy development and deployment, RoboDojo enables policies to be integrated, evaluated, and compared under a shared protocol.

Our experiments show that current robot policies remain far from reliable general-purpose manipulation. In simulation, RoboDojo reveals insufficient capability coverage across key dimensions, including generalization, long-horizon execution, precise manipulation, memory, and open-ended task understanding, indicating that current models are not yet comprehensive across diverse manipulation requirements. In the real world, the evaluated policies still perform poorly on complex manipulation tasks, showing that current models are not yet robust enough to handle challenging physical deployment conditions. RoboDojo will continue to update its leaderboard with new policies, tasks, evaluation results, and community submissions, providing a challenging testbed, reproducible evaluation platform, and open benchmark for tracking progress toward more capable omni-manipulation policies.
